% META: {"id": "1NSI-TC-EX-026", "chapitre": "1NSI-TYPES-CONSTRUITS", "type_objet": "exercice", "capacites": ["1NSI-TYPES-CONSTRUITS-C3"], "parcours": 1, "competences": ["analyser"], "duree_min": 8, "mode_creation": "ex_nihilo", "sources_inspiration": [], "status": "needs_review"}
% BEGIN-TRACE
% carte = [[20, 22, 19],
%           [25, 28, 23],
%           [18, 17, 21]]
% total = 0
% for ligne in carte:
%     for val in ligne:
%         total = total + val
% print(total)
% EXPECTED
% 193
% END-TRACE
\begin{exercice}{1NSI-TC-EX-026}{1}{8}
Une grille représente les températures relevées en différents points d'une carte :
\begin{python}
carte = [[20, 22, 19],
          [25, 28, 23],
          [18, 17, 21]]
total = 0
for ligne in carte:
    for val in ligne:
        total = total + val
print(total)
\end{python}
\begin{enumerate}
  \item Donner, sans exécuter, l'affichage produit.
  \item Modifier le code pour qu'il affiche la moyenne des températures de la grille.
\end{enumerate}
\end{exercice}
